% Chapter 14: American options by simulation.

An American option can be exercised at any time up to its maturity; a Bermudan option, on a
finite set of dates. Pricing either means solving an optimal stopping problem under the
risk-neutral measure: at each date the holder compares the payoff of exercising now with
the value of keeping the option alive. In the binomial model of the IEOR 222 notes this is
a backward induction over the nodes of a tree, and the trees and finite differences of
\cref{ch:20-grids} do the same on a grid. Grids are efficient in one or two state
variables and hopeless in ten (\cref{prop:mc-dimension}). Monte Carlo handles many state
variables, but it runs forward in time, while the exercise decision at a date requires the
conditional expectation of the future, which one simulated path does not reveal.

Three ideas make simulation work, and each is a section of this chapter.
\Cref{sec:am-dp} reduces the problem on a finite set of dates to a dynamic programming
recursion whose only unknowns are the continuation values, conditional expectations of a
function of the next state. \Cref{sec:am-lsm} estimates them by regression across
simulated paths, as \textcite{longstaff-schwartz} proposed: this produces an exercise
policy, and any policy gives a \emph{lower} bound for the price. \Cref{sec:am-duality}
turns the problem around: every martingale gives an \emph{upper} bound, the optimal one
gives the price exactly, and a good policy gives a good martingale
\citep{rogers-duality,haugh-kogan,andersen-broadie}. The two bounds bracket the price and
measure how good the policy is. The words ``low'' and ``high'' need care, and we will be
precise about which estimator is biased in which direction.

Throughout, $S$ is the Black--Scholes stock with yield $y$ under the risk-neutral measure
$\Q$ (\cref{def:sm-model}), with a filtration $(\F_t)$ for its Brownian motion; all
expectations are under $\Q$, and we write $\E$ for $\EQ$. The running examples are the put of
\textcite{longstaff-schwartz}, with $K=40$, $r=6\%$, $y=0$, $\sigma\in\{20\%,40\%\}$,
$T\in\{1,2\}$, $S_0\in\{36,38,\dots,44\}$ and $50$ exercise dates per year, and the
course's test case $S_0=K=100$, $T=1$, $r=1\%$, $y=2\%$, $\sigma=20\%$.

\section{Bermudan options and optimal stopping}\label{sec:am-dp}

Fix exercise dates $0=t_0<t_1<\dots<t_m=T$ and write $\F_k:=\F_{t_k}$. A Bermudan claim
with payoff function $g$ pays $g(S_{t_k})$ if it is exercised at $t_k$; let
\[
  Z_k=e^{-rt_k}g(S_{t_k}),\qquad k=0,\dots,m,
\]
be the payoff discounted to time $0$. More generally, $Z$ may be any adapted sequence. A
\emph{stopping time} is a random variable $\tau$ with values in $\{0,\dots,m\}$ such that
$\{\tau=j\}\in\F_j$ for every $j$: the decision to stop at $t_j$ uses only the information
available at $t_j$. Write $\mathcal T_k$ for the set of stopping times with $k\le\tau\le m$.
The price of the Bermudan claim is
\begin{equation}\label{eq:am-price}
  U_0=\sup_{\tau\in\mathcal T_0}\E[Z_\tau],
\end{equation}
the value of the holder's optimal stopping problem. That this is the arbitrage-free price,
the smallest initial capital from which a self-financing portfolio dominates the payoff at
every date, is shown for finite probability spaces in the IEOR 222 notes and in general
discrete time by \textcite{follmer-schied}, Chapter~6; \textcite{shreve2}, Section~8.4,
uses the corresponding definition for the American put. Here we take
\eqref{eq:am-price} as the definition.

On a finite set of dates, integrability of $\max_k|Z_k|$ is the same as integrability of
each $Z_k$, because $\max_k|Z_k|\le\sum_k|Z_k|$; for calls and puts it holds since
$\E[S_t]<\infty$. The next lemma is optional stopping on a finite horizon, conditionally on
$\F_k$.

\begin{lemma}[Optional stopping]\label{lem:am-optional}
Let $(Y_j)_{k\le j\le m}$ be a supermartingale with respect to $(\F_j)$ and
$\tau\in\mathcal T_k$. Then $Y_\tau$ is integrable and $\E[Y_\tau\mid\F_k]\le Y_k$, with
equality if $Y$ is a martingale.
\end{lemma}

\begin{proof}
$|Y_\tau|\le\sum_j|Y_j|$ is integrable, and $Y_\tau-Y_k=\sum_{j=k+1}^m\1_{\{\tau\ge
j\}}(Y_j-Y_{j-1})$. The event $\{\tau\ge j\}=\Omega\setminus\{\tau\le j-1\}$ is in
$\F_{j-1}$, so by the tower property and taking out what is known,
$\E[\1_{\{\tau\ge j\}}(Y_j-Y_{j-1})\mid\F_k]=\E\big[\1_{\{\tau\ge
j\}}\E[Y_j-Y_{j-1}\mid\F_{j-1}]\,\big|\,\F_k\big]\le0$, with equality for a martingale.
\end{proof}

\begin{theorem}[Bermudan dynamic programming]\label{thm:am-dp}
Let $Z$ be adapted with $\E[\max_k|Z_k|]<\infty$. Define $U_m=Z_m$ and, for
$k=m-1,\dots,0$,
\begin{equation}\label{eq:am-dp}
  C_k=\E[U_{k+1}\mid\F_k],\qquad U_k=\max(Z_k,C_k).
\end{equation}
\begin{enumerate}[label=(\roman*)]
  \item (\emph{Snell envelope}.) Each $U_k$ is integrable, $U\ge Z$, and $U$ is a
    supermartingale. Every supermartingale $V$ with $V\ge Z$ satisfies $V\ge U$.
  \item (\emph{Optimal stopping}.) For every $k$ and every $\tau\in\mathcal T_k$,
    $\E[Z_\tau\mid\F_k]\le U_k$, with equality for
    $\tau^*_k=\min\{j\ge k:\ j=m\text{ or }Z_j\ge C_j\}$. In particular
    $U_0=\max_{\tau\in\mathcal T_0}\E[Z_\tau]=\E[Z_{\tau^*_0}]$ is the price \eqref{eq:am-price}.
  \item (\emph{A sufficient condition}.) If $\tau\in\mathcal T_k$, $Z_\tau=U_\tau$, and the
    stopped process $(U_{j\wedge\tau})_{k\le j\le m}$ is a martingale, then
    $\E[Z_\tau\mid\F_k]=U_k$: $\tau$ is optimal from $t_k$ on.
  \item (\emph{Stopping-time form}.) $\tau^*_m=m$, $\tau^*_k=k$ on $\{Z_k\ge C_k\}$ and
    $\tau^*_k=\tau^*_{k+1}$ elsewhere; and $C_k=\E[Z_{\tau^*_{k+1}}\mid\F_k]$.
  \item (\emph{Doob decomposition}.) Let $M^*_0=A^*_0=0$ and
    \begin{equation}\label{eq:am-doob}
      M^*_{k+1}-M^*_k=U_{k+1}-C_k,\qquad A^*_{k+1}-A^*_k=U_k-C_k .
    \end{equation}
    Then $M^*$ is a martingale, $A^*$ is nondecreasing and predictable ($A^*_{k+1}$ is
    $\F_k$-measurable), and $U_k=U_0+M^*_k-A^*_k$.
  \item (\emph{Markov form}.) If $Z_k=e^{-rt_k}g(S_{t_k})$ with $g$ measurable and
    $|g(s)|\le a(1+s)$, then $U_k=e^{-rt_k}V_k(S_{t_k})$ and $C_k=e^{-rt_k}c_k(S_{t_k})$,
    where $V_m=g$ and, with $\Delta_k=t_{k+1}-t_k$ and $\xi\sim\mathcal N(0,1)$,
    \begin{equation}\label{eq:am-markov}
      c_k(s)=e^{-r\Delta_k}\,\E\big[V_{k+1}\big(se^{\nu\Delta_k+\sigma\sqrt{\Delta_k}\,\xi}\big)\big],
      \qquad V_k=\max(g,c_k).
    \end{equation}
\end{enumerate}
\end{theorem}

\begin{proof}
(i) Let $Z^\star=\max_j|Z_j|$. By backward induction $|U_k|\le\E[Z^\star\mid\F_k]$: this
holds for $k=m$, and if it holds for $k+1$ then $|C_k|\le\E[|U_{k+1}|\mid\F_k]\le
\E[Z^\star\mid\F_k]$ and $|Z_k|\le\E[Z^\star\mid\F_k]$ because $Z_k$ is $\F_k$-measurable.
So $U$ is integrable and adapted, $U\ge Z$, and $U_k\ge C_k=\E[U_{k+1}\mid\F_k]$ is the
supermartingale property. If $V\ge Z$ is a supermartingale, then $V_m\ge Z_m=U_m$, and
$V_{k+1}\ge U_{k+1}$ implies $V_k\ge\E[V_{k+1}\mid\F_k]\ge C_k$ and $V_k\ge Z_k$, so
$V_k\ge U_k$.

(iii) The stopped process is a martingale, so
$U_k=\E[U_{m\wedge\tau}\mid\F_k]=\E[U_\tau\mid\F_k]=\E[Z_\tau\mid\F_k]$.

(ii) By \cref{lem:am-optional} and (i), $U_k\ge\E[U_\tau\mid\F_k]\ge\E[Z_\tau\mid\F_k]$.
For $\tau^*_k$: if $\tau^*_k=j<m$ then $Z_j\ge C_j$, so $U_j=Z_j$; and $U_m=Z_m$. The
increments of the stopped process are
$U_{(j+1)\wedge\tau^*_k}-U_{j\wedge\tau^*_k}=\1_{\{\tau^*_k>j\}}(U_{j+1}-U_j)$, and on
$\{\tau^*_k>j\}\in\F_j$ we have $Z_j<C_j$, hence $U_j=C_j$ and
$\E[\1_{\{\tau^*_k>j\}}(U_{j+1}-U_j)\mid\F_j]=\1_{\{\tau^*_k>j\}}(C_j-U_j)=0$. So (iii)
applies. As $\F_0$ is trivial, the maximum over $\mathcal T_0$ is attained at $\tau^*_0$.

(iv) The recursion is the definition of $\tau^*_k$ read from date $k$. By (ii) at $k+1$
and the tower property, $C_k=\E\big[\E[Z_{\tau^*_{k+1}}\mid\F_{k+1}]\mid\F_k\big]=\E[Z_{\tau^*_{k+1}}\mid\F_k]$.

(v) $\E[U_{k+1}-C_k\mid\F_k]=0$, $U_k-C_k\ge0$ is $\F_k$-measurable, and the sum of the two
increments in \eqref{eq:am-doob} is $U_{k+1}-U_k$.

(vi) By backward induction. Suppose $U_{k+1}=e^{-rt_{k+1}}V_{k+1}(S_{t_{k+1}})$ with
$|V_{k+1}(s)|\le a'(1+s)$. Write $S_{t_{k+1}}=S_{t_k}e^{\nu\Delta_k+\sigma(W_{t_{k+1}}-W_{t_k})}$;
the increment of $W$ is independent of $\F_k$ and $S_{t_k}$ is $\F_k$-measurable, so the
independence lemma \citep[Lemma~2.3.4]{shreve2} gives $C_k=e^{-rt_k}c_k(S_{t_k})$ with
$c_k$ as in \eqref{eq:am-markov}, which is finite with $|c_k(s)|\le
a'(1+se^{(r-y)\Delta_k})e^{-r\Delta_k}$. Then $U_k=e^{-rt_k}\max(g,c_k)(S_{t_k})$, and
$V_k$ again grows at most linearly.
\end{proof}

The process $U$ is the \emph{Snell envelope} of $Z$; in the Markov form, $V_k$ is the value
of the option at $t_k$ and $c_k$ its \emph{continuation value}, the value of holding it for
one more period. Part (iv) is the form of the recursion that simulation will use: it moves
stopping times, not values, backwards, and the continuation value is the conditional mean
of the discounted payoff received by following the optimal rule from the next date on. On
continuous time, with its additional measurability issues, see \textcite{peskir-shiryaev},
Chapter~I, and for the financial formulation \textcite{glasserman-mc}, Section~8.1.

\begin{proposition}[From Bermudan to American]\label{prop:am-american}
Let $g$ be continuous with $|g(s)|\le a(1+s)$, $Z_t=e^{-rt}g(S_t)$ for $t\in[0,T]$, and
let the American price be $A=\sup_\tau\E[Z_\tau]$ over stopping times $\tau$ of $(\F_t)$
with values in $[0,T]$. For a set of dates $\Pi$ as above, let $B(\Pi)$ be the Bermudan price
\eqref{eq:am-price}. Then $B(\Pi)\le B(\Pi')\le A$ whenever $\Pi\subseteq\Pi'$, and
$B(\Pi_n)\to A$ for every sequence $(\Pi_n)$ whose mesh tends to $0$.
\end{proposition}

\begin{proof}
A stopping time with values in $\Pi$ is one with values in $\Pi'$, and one with values in
$[0,T]$: the sets over which the suprema are taken increase. For the limit, $\E[\sup_{t\le
T}|Z_t|]<\infty$: $|Z_t|\le a(1+S_t)$ and $S_t=S_0e^{(r-y)t}\mathcal E_t\le S_0e^{|r-y|T}\mathcal E_t$
with the martingale $\mathcal E_t=e^{\sigma W_t-\sigma^2t/2}$, whose supremum over $[0,T]$ is
square-integrable by Doob's maximal inequality \citep[Theorem~1.3.8]{karatzas-shreve}. Fix a stopping
time $\tau$ and let $\tau_n=\min\{t\in\Pi_n:t\ge\tau\}$. Since
$\{\tau_n\le t\}=\{\tau\le t\}\in\F_t$ for $t\in\Pi_n$, $\tau_n$ is a stopping time with
values in $\Pi_n$, so $B(\Pi_n)\ge\E[Z_{\tau_n}]$. As the mesh tends to $0$, $\tau_n\to\tau$,
so $Z_{\tau_n}\to Z_\tau$ by continuity of the paths, and dominated convergence (with the
bound $\sup_t|Z_t|$) gives $\E[Z_{\tau_n}]\to\E[Z_\tau]$. Hence $\liminf_nB(\Pi_n)\ge\E[Z_\tau]$
for every $\tau$, so $\liminf_nB(\Pi_n)\ge A$.
\end{proof}

The proposition gives no rate. On the Longstaff--Schwartz put with $S_0=40$, $\sigma=20\%$,
$T=1$, our grid (\cref{rem:am-references}) gives $2.0664$ (European), $2.1991$, $2.2930$,
$2.3141$, $2.3185$ for $m=1,2,10,50,250$ equally spaced dates, and the American price is
$2.3196$: the gap falls from $0.027$ to $0.0055$ to $0.0011$ as $m$ goes from $10$ to $50$
to $250$, roughly like $1/m$ in this example.

\begin{remark}[Reference values]\label{rem:am-references}
Our reference prices use no simulation (\texttt{bermudan\_reference} and
\texttt{american\_binomial} in the code of this chapter). Bermudan prices come from \eqref{eq:am-markov} on a grid in $\log s$ of mesh $h$: the trapezoid rule
against the exact Gaussian transition density (an FFT convolution), extrapolated in $h$; the
error is below $10^{-6}$ here, and a binomial tree that exercises only at the $50$ dates
agrees. American prices come from the binomial tree with the Black--Scholes value in its last
step and Richardson extrapolation \citep{broadie-detemple}: with $2000$ steps in
\cref{tab:am-ls}, quoted to $10^{-3}$, and $4000$ steps elsewhere, with an error of about
$10^{-5}$. Grids are the subject of \cref{ch:20-grids}.
\end{remark}

\begin{remark}[Degenerate cases]\label{rem:am-degenerate}
\begin{enumerate}[label=(\alph*)]
  \item \emph{Calls without dividends are never exercised early.} If $y=0$, $r\ge0$ and
    $K\ge0$, Jensen's inequality and the martingale property of $e^{-rt}S_t$ give
    $\E[Z_{k+1}\mid\F_k]\ge\big(e^{-rt_k}S_{t_k}-e^{-rt_{k+1}}K\big)^+\ge Z_k$ for
    $Z_k=e^{-rt_k}(S_{t_k}-K)^+$: $Z$ is a submartingale. So $V_k=\E[Z_m\mid\F_k]$ is a
    martingale with $V\ge Z$, and \cref{thm:am-dp}(i) gives $U_0\le V_0=\E[Z_m]\le U_0$. The
    Bermudan call, and by \cref{prop:am-american} the American call, is worth the European
    call. With $y>0$ this fails: on the test case, where $y>r$, the Bermudan call with $50$
    dates is worth $7.4567$ and the American call $7.4597$, against $7.36429$ for the
    European call.
  \item \emph{The test-case put.} It is worth $8.349406$ as a European option, $8.349408$ as a
    Bermudan option with $50$ dates and as an American option: early exercise is worth
    $2\times10^{-6}$. The reason is where the exercise region lies. Since
    $c_k(s)\ge e^{-r\Delta}\E[K-S_{t_{k+1}}\mid S_{t_k}=s]=Ke^{-r\Delta}-se^{-y\Delta}$ with
    $\Delta=t_{k+1}-t_k$, exercise at $t_k$ is suboptimal whenever
    $s>\bar s=K(1-e^{-r\Delta})/(1-e^{-y\Delta})$, which is close to $Kr/y=50$ here. The
    boundary rises from $44.9$ at $t_1$ to $\bar s=50.005$ just before maturity, prices that
    $S$ rarely reaches within a year. A put deep enough in the money is still exercised.
  \item \emph{Exercise at $t_0$.} $U_0=\max(Z_0,C_0)$, and immediate exercise is optimal
    when $Z_0\ge C_0$, as for a put deep in the money. All simulated paths share $S_{t_0}$, so
    no regression is possible at $t_0$: $C_0$ is estimated by a sample mean
    (\exref{ex:am-exercise-now}).
  \item \emph{Out of the money.} If $g\ge0$ then $C_k\ge0$, so stopping where $Z_k=0$ is
    never better than continuing, and a rule may as well continue there.
\end{enumerate}
\end{remark}

\section{Regression: the Longstaff--Schwartz algorithm}\label{sec:am-lsm}

The recursion of \cref{thm:am-dp}(iv) needs the continuation value $C_k$ at the simulated
state. Estimating it by a fresh simulation from every node of every path, as in the random
tree of \textcite{broadie-glasserman-tree}, costs a number of paths exponential in $m$.
The idea of \textcite{carriere}, \textcite{tsitsiklis-vanroy} and
\textcite{longstaff-schwartz} is to estimate the function $c_k$ in one go, by regressing
across paths: since $C_k=\E[Z_{\tau^*_{k+1}}\mid S_{t_k}]$, the realised discounted cash
flows of the paths are noisy observations of it, and least squares onto functions of
$S_{t_k}$ recovers their conditional mean. Let $\psi=(\psi_1,\dots,\psi_J)$ be a vector of
basis functions.

\paragraph{The algorithm.} Simulate $N$ training paths $S^{(i)}$ at the dates $t_k$ exactly
(\cref{prop:ps-gbm-exact}), with discounted payoffs $Z^{(i)}_k$, and set
$Y^{(i)}=Z^{(i)}_m$. For $k=m-1,\dots,1$:
\begin{enumerate}[label=(\arabic*)]
  \item let $I_k=\{i:Z^{(i)}_k>0\}$ be the paths in the money at $t_k$;
  \item compute $\hat\beta_k\in\argmin_\beta\sum_{i\in I_k}\big(Y^{(i)}-\beta^\top\psi(S^{(i)}_{t_k})\big)^2$
    and set $\hat C_k(s)=\hat\beta_k^\top\psi(s)$;
  \item for $i\in I_k$ with $Z^{(i)}_k\ge\hat C_k(S^{(i)}_{t_k})$, replace $Y^{(i)}$ by
    $Z^{(i)}_k$.
\end{enumerate}
Finally $\hat c_0=\frac1N\sum_iY^{(i)}$, and the \emph{in-sample estimate} is
$\hat V^{\rm in}_N=\max(Z_0,\hat c_0)$. The coefficients define an exercise rule: for a
path $S$,
\begin{equation}\label{eq:am-rule}
  \hat\tau=\begin{cases}0&\text{if }Z_0>0\text{ and }Z_0\ge\hat c_0,\\
  \min\big\{k\ge1:\ k=m\text{ or }\big(Z_k>0\text{ and }Z_k\ge\hat C_k(S_{t_k})\big)\big\}&\text{otherwise.}\end{cases}
\end{equation}
The \emph{out-of-sample estimate} freezes $\hat\beta_1,\dots,\hat\beta_{m-1}$ and $\hat c_0$,
simulates $n$ new paths independent of the training ones, and averages
$Z^{(i)}_{\hat\tau^{(i)}}$: $\hat V^{\rm out}_n=\frac1n\sum_{i=1}^nZ^{(i)}_{\hat\tau^{(i)}}$.

Step~(3) is \cref{thm:am-dp}(iv) with $\hat C_k$ in place of $C_k$: after it, $Y^{(i)}$ is the
discounted cash flow of path $i$ under the estimated rule from $t_k$ on, and the regression
at $t_{k-1}$ uses these cash flows, not the fitted values. Step~(2) is the least-squares
regression of \cref{thm:cv-multiple}(iii), now used to estimate a conditional expectation
rather than an unconditional one. Regressing on in-the-money paths only is the choice of
\textcite{longstaff-schwartz}: the decision matters only there (\cref{rem:am-degenerate}(d)),
and the fit is not spent on the region $S\ge K$ where the put is out of the money.

\begin{incode}
\begin{lstlisting}
Y = D[m] * g(S[:, m])                     # discounted cash flows, time-0 units
for k in range(m - 1, 0, -1):
    Zk = D[k] * g(S[:, k]);  idx = np.flatnonzero(Zk > 0)
    X = basis(S[idx, k])                  # scaled basis, e.g. x = S / K
    beta, *_ = np.linalg.lstsq(X, Y[idx], rcond=None)
    ex = Zk[idx] >= X @ beta;  Y[idx[ex]] = Zk[idx[ex]]
\end{lstlisting}
\texttt{lsm\_fit}, \texttt{policy\_cashflows} (the frozen rule) and \texttt{lsm\_price} (both
estimates, and a control variate) in \coderef{code/src/mfe230d/american.py}.
\end{incode}

\begin{proposition}[Longstaff--Schwartz: which way the estimates are biased]\label{prop:am-lsm}
Let $g\ge0$, $\E[Z_k]<\infty$ for every $k$, and let $\mathcal G$ be the $\sigma$-algebra
generated by the training paths.
\begin{enumerate}[label=(\roman*)]
  \item (\emph{Out of sample: low.}) Let the evaluation paths each have the law of $S$, and
    let them, as a family, be independent of $\mathcal G$. For every value of the training
    data, the rule \eqref{eq:am-rule} is a stopping time in $\mathcal T_0$. Its value
    $\hat L=\E[Z_{\hat\tau}\mid\mathcal G]$ satisfies $\hat L\le U_0$, the out-of-sample
    estimator is conditionally unbiased for it, $\E[\hat V^{\rm out}_n\mid\mathcal G]=\hat L$,
    and hence
    \[
      \E\big[\hat V^{\rm out}_n\big]\le U_0 .
    \]
    If the evaluation paths are i.i.d.\ (or i.i.d.\ antithetic pairs, whose averages are then
    the samples), then $\hat V^{\rm out}_n\to\hat L$ almost surely as $n\to\infty$, and when
    $\E[Z_k^2]<\infty$ the central limit theorem (\cref{thm:mc-clt}) applies conditionally
    on $\mathcal G$. If
    $\hat C_k(S_{t_k})=C_k$ almost surely on $\{Z_k>0\}$ for $1\le k\le m-1$ and
    $\{Z_0>0,Z_0\ge\hat c_0\}=\{Z_0>0,Z_0\ge C_0\}$, then $\hat L=U_0$.
  \item (\emph{In sample: no guarantee.}) Suppose that at every date the matrix
    $\big(\psi_j(S^{(i)}_{t_k})\big)_{i\in I_k,\,j\le J}$ has rank $|I_k|$, so that least
    squares interpolates the cash flows. Then after the backward pass
    $Y^{(i)}=\max_{1\le k\le m}Z^{(i)}_k$ for every $i$, and
    \[
      \E\big[\hat V^{\rm in}_N\big]\ \ge\ \max\big(Z_0,\E[\max_{1\le k\le m}Z_k]\big)\ \ge\ U_0 .
    \]
    The last inequality is an equality if $Z_0\ge\E[\max_{k\ge1}Z_k]$; otherwise it is strict
    unless the optimal rule from $t_1$ on stops, almost surely, at a date where the path
    attains $\max_{k\ge1}Z_k$.
\end{enumerate}
\end{proposition}

\begin{proof}
(i) Fix the training data, so that $\hat\beta_k$ and $\hat c_0$ are numbers. Then
$\{\hat\tau=k\}$ is determined by $S_{t_0},\dots,S_{t_k}$, so $\hat\tau\in\mathcal T_0$, and
$\E[Z_{\hat\tau}]\le U_0$ by \cref{thm:am-dp}(ii). The evaluation paths are independent of
$\mathcal G$, so by the independence lemma $\hat L$ is this expectation evaluated at the
training data, hence $\hat L\le U_0$; the same argument for each evaluation path gives
$\E[\hat V^{\rm out}_n\mid\mathcal G]=\hat L$, and taking expectations gives the inequality.
For i.i.d.\ paths (or pairs) the terms are conditionally i.i.d.\ given $\mathcal G$, so the
strong law and the central limit theorem apply conditionally. For the equality case, use \cref{thm:am-dp}(iii): at
$\hat\tau$ we have $Z=U$ (where the rule stops before $m$, $Z_k\ge C_k$); and on
$\{\hat\tau>j\}$ either $Z_j>0$ and $Z_j<C_j$, or $Z_j=0\le C_j$ (as $g\ge0$), and in both
cases $U_j=C_j$, so the stopped Snell envelope is a martingale as in the proof of
\cref{thm:am-dp}(ii).

(ii) Rank $|I_k|$ means the fitted values can match any vector of responses, so the
least-squares residuals vanish: $\hat C_k(S^{(i)}_{t_k})$ is the current $Y^{(i)}$ for
$i\in I_k$. Step~(3) then replaces $Y^{(i)}$ by $Z^{(i)}_k$ exactly when $Z^{(i)}_k\ge Y^{(i)}$,
that is, $Y^{(i)}\leftarrow\max(Z^{(i)}_k,Y^{(i)})$; for $i\notin I_k$, $Z^{(i)}_k=0\le
Y^{(i)}$ and the same formula holds. By induction, $Y^{(i)}=\max_{j\ge k}Z^{(i)}_j$ after
date $k$. The map $x\mapsto\max(Z_0,x)$ is convex, so by Jensen's inequality
$\E[\hat V^{\rm in}_N]\ge\max(Z_0,\E[\max_{k\ge1}Z_k])$. Finally
$\E[\max_{k\ge1}Z_k]\ge\E[Z_{\tau^*_1}]=C_0$ by \cref{thm:am-dp}(iv), and
$\max(Z_0,C_0)=U_0$. If $Z_0<\E[\max_{k\ge1}Z_k]$, the two maxima differ unless
$C_0=\E[\max_{k\ge1}Z_k]$, that is, unless $\max_{k\ge1}Z_k=Z_{\tau^*_1}$ almost surely.
\end{proof}

Part~(ii) is an extreme case (our code skips the regression at dates with at most $J$
in-the-money paths, so it never interpolates), but it shows the mechanism. In-sample, each
path's own future enters the fit that decides when that path stops: a path whose future
cash flow is high pulls its fitted continuation value up and is kept alive, and one whose
future is low is exercised. That is foresight, and it biases upwards. Out of sample, the
fitted function is just some function, the rule just some stopping time, and its value can
only be below the optimum. \Cref{fig:am-bias} shows both effects on the
Longstaff--Schwartz put with $S_0=40$, $\sigma=20\%$, $T=1$. With $N=100$ training paths
and the basis $1,x,x^2$ ($J=3$, $x=S/K$), the in-sample estimate is on average $0.28$
above the price $2.3141$, while the value of the fitted policy is on average $0.157$
below it; each replication gives a policy worth less than the price, as (i) requires. The
in-sample bias shrinks as $N$ grows; with $J=3$ it is $-0.011$ at $N=10^4$, with a standard
error of $0.005$, as the in-sample estimate approaches the value of the limiting policy, which
is below the price.
More basis functions make the in-sample estimate worse for small $N$, and the policy better
only once $N$ is large.
\textcite{longstaff-schwartz} found in-sample and out-of-sample estimates ``virtually
identical'' with $10^5$ paths and four basis functions, which \cref{tab:am-ls} confirms;
with few paths or many basis functions only the out-of-sample estimate is a lower bound.

\begin{figure}[ht]
  \centering
  \includegraphics[width=0.82\textwidth]{ch14-bias}
  \caption{Longstaff--Schwartz put, $S_0=K=40$, $\sigma=20\%$, $T=1$, $50$ dates (price
    $2.3141$). For each number $N$ of training paths, LSM is fitted on $150$, $80$, $40$,
    $30$ and $20$ independent sets of $N$ paths (for $N=100,300,1000,3000,10^4$) with the
    monomials $1,x,\dots,x^{J-1}$ in $x=S/K$, $J=3$ and $J=9$. Solid: mean in-sample
    estimate minus the price. Dashed: mean value of the fitted policy minus the price, the
    value computed exactly on the grid, which is the conditional mean of the out-of-sample
    estimator given the training paths. Bars: two standard errors across the sets. Our
    implementation (NumPy \texttt{default\_rng}, seeds $[8,J-1,N,j]$).}
  \label{fig:am-bias}
\end{figure}

\begin{remark}[Convergence]\label{rem:am-clp}
\textcite{clement-lamberton-protter} analyse a variant that regresses on all paths, without
the restriction $Z_k>0$, and exercises where $Z_k$ is at least the fitted value; it is not
the rule \eqref{eq:am-rule}. Let $X_k=S_{t_k}$, $Z_k=f_k(X_k)$ with $\E[Z_k^2]<\infty$, and let
$\psi_1,\psi_2,\dots$ be functions such that, at every date $1\le k\le m-1$,
$\psi_1(X_k),\dots,\psi_J(X_k)$ are linearly independent in $L^2$ for every $J$. Let
$\tau^{[J]}$ be their rule built from the exact $L^2$ projections onto the first $J$
functions (the regression with infinitely many paths). Their Theorem~3.1: if moreover the
whole sequence is total in $L^2(\sigma(X_k))$, then as $J\to\infty$,
$\E[Z_{\tau^{[J]}_k}\mid\F_k]\to\E[Z_{\tau^*_k}\mid\F_k]$ in $L^2$ for every $k$, so the value
of $\tau^{[J]}$ tends to $U_0$. Their Theorem~3.2, which needs linear independence but not
totality: for fixed $J$, if the projected continuation value ties with the payoff with
probability zero at every date, the in-sample estimate converges almost surely to the value
of $\tau^{[J]}$ as $N\to\infty$, which is at most $U_0$ by \cref{thm:am-dp}(ii); they also
prove a central limit theorem for this convergence. The two limits are taken one after the other. How fast $N$ must grow with $J$
is studied by \textcite{glasserman-yu}: for polynomial bases, the number of paths needed
for convergence in the worst case grows exponentially in the degree for Brownian motion, and
faster for geometric Brownian motion. See also \textcite{glasserman-mc}, Section~8.6.
\end{remark}

\begin{remark}[Bases and conditioning]\label{rem:am-basis}
The basis should be scaled. \textcite{longstaff-schwartz} use a constant and the weighted
Laguerre functions $e^{-x/2}$, $e^{-x/2}(1-x)$ and $e^{-x/2}(1-2x+\frac{x^2}2)$; we use them
in the scaled variable $x=S/K$, because unscaled columns of very different sizes make the
regression matrix badly conditioned. On the
in-the-money prices at $t=0.5$ of the put with $S_0=36$, $\sigma=20\%$ ($10^5$ paths, seed
$0$), the condition number of the regression matrix is $2.6\times10^7$ for $1,S,S^2,S^3$ and
$5.9\times10^3$ for $1,x,x^2,x^3$; $1.1\times10^8$ and $1.4\times10^5$ for the Laguerre basis
in $S$ and in $x$; and $2.7\times10^9$ for the monomials of degree $8$ in $x$. Solving the
normal equations squares these numbers; \texttt{numpy.linalg.lstsq} works with the matrix
itself (by a singular value decomposition) and returns a least-squares solution even when
the matrix is rank-deficient. Our code skips the regression, and never exercises, at a date
with at most $J$ usable paths: any rule that depends only on the path up to $t_k$ is a
stopping time, so \cref{prop:am-lsm}(i) is not affected.
\end{remark}

\paragraph{Table 1 of Longstaff and Schwartz.} \Cref{tab:am-ls} reproduces the table with
which \textcite{longstaff-schwartz} introduced the method, with their setup: $10^5$ paths
($5\times10^4$ and their antithetic images), $50$ dates per year, the four basis functions
above. Our closed-form European prices agree with theirs to the three printed decimals. Two
findings about their columns. First, their ``finite-difference American'' prices agree with
our \emph{Bermudan} prices to within $0.0005$ in the $15$ rows with $\sigma=20\%$ or $T=1$,
and lie between our Bermudan and American prices, up to the last printed digit, in the five
rows with $\sigma=40\%$, $T=2$; the American prices exceed the Bermudan ones by $0.003$ to
$0.009$. They do not say whether their grid allowed exercise at every time step; the
comparison suggests that their column approximates the Bermudan price, which is also what
LSM estimates. Second, their standard errors are close to the naive formula $\hat s_n/\sqrt
n$ over all $10^5$ paths (we get $0.0092$, $0.0086$, $0.0206$ for the first, ninth and last
rows; they print $0.010$, $0.009$, $0.021$). But antithetic pairs are negatively correlated:
the correlation $\rho$ of the two cash flows of a pair is between $-0.68$ and $-0.56$ in these
rows, and the standard error from the $5\times10^4$ pair averages
(\cref{prop:cv-antithetic}), which we report, is $\sqrt{1+\rho}$ times the naive one, that
is $0.57$ to $0.66$ times. All our in-sample estimates are within $1.3$
of our standard errors of the Bermudan price, and the out-of-sample estimates within $2.4$,
with mean deviations of $-0.2$ and $-0.3$ standard errors: the bias is small compared with
the noise at this sample size.

\begin{table}[ht]
  \centering
  \footnotesize
  \setlength{\tabcolsep}{4.2pt}
  \begin{tabular}{@{}ccc@{\hspace{8pt}}cc@{\hspace{8pt}}cc@{\hspace{8pt}}cc@{}}
    \toprule
    & & & \multicolumn{2}{c}{ours, no simulation} & \multicolumn{2}{c}{Longstaff--Schwartz} & \multicolumn{2}{c}{ours, LSM (s.e.)}\\
    \cmidrule(lr){4-5}\cmidrule(lr){6-7}\cmidrule(l){8-9}
    $S_0$ & $\sigma$ & $T$ & Bermudan & American & ``FD'' & LSM (s.e.) & in-sample & out-of-sample\\
    \midrule
36 & 0.2 & 1 & 4.478 & 4.487 & 4.478 & 4.472 (.010) & 4.482 (.006) & 4.470 (.006)\\
36 & 0.2 & 2 & 4.840 & 4.848 & 4.840 & 4.821 (.012) & 4.839 (.007) & 4.841 (.007)\\
36 & 0.4 & 1 & 7.101 & 7.109 & 7.101 & 7.091 (.020) & 7.104 (.008) & 7.088 (.008)\\
36 & 0.4 & 2 & 8.507 & 8.514 & 8.508 & 8.488 (.024) & 8.517 (.010) & 8.507 (.010)\\
38 & 0.2 & 1 & 3.250 & 3.257 & 3.250 & 3.244 (.009) & 3.247 (.005) & 3.253 (.005)\\
38 & 0.2 & 2 & 3.745 & 3.751 & 3.745 & 3.735 (.011) & 3.745 (.006) & 3.751 (.006)\\
38 & 0.4 & 1 & 6.148 & 6.155 & 6.148 & 6.139 (.019) & 6.138 (.008) & 6.145 (.008)\\
38 & 0.4 & 2 & 7.668 & 7.675 & 7.670 & 7.669 (.022) & 7.675 (.010) & 7.674 (.010)\\
40 & 0.2 & 1 & 2.314 & 2.320 & 2.314 & 2.313 (.009) & 2.309 (.005) & 2.310 (.005)\\
40 & 0.2 & 2 & 2.885 & 2.890 & 2.885 & 2.879 (.010) & 2.879 (.006) & 2.877 (.006)\\
40 & 0.4 & 1 & 5.312 & 5.318 & 5.312 & 5.308 (.018) & 5.309 (.009) & 5.305 (.009)\\
40 & 0.4 & 2 & 6.917 & 6.923 & 6.920 & 6.921 (.022) & 6.906 (.010) & 6.919 (.010)\\
42 & 0.2 & 1 & 1.617 & 1.621 & 1.617 & 1.617 (.007) & 1.624 (.006) & 1.607 (.006)\\
42 & 0.2 & 2 & 2.212 & 2.217 & 2.212 & 2.206 (.010) & 2.207 (.007) & 2.219 (.007)\\
42 & 0.4 & 1 & 4.582 & 4.588 & 4.582 & 4.588 (.017) & 4.578 (.010) & 4.559 (.010)\\
42 & 0.4 & 2 & 6.244 & 6.250 & 6.248 & 6.243 (.021) & 6.231 (.011) & 6.254 (.011)\\
44 & 0.2 & 1 & 1.110 & 1.113 & 1.110 & 1.118 (.007) & 1.111 (.006) & 1.112 (.006)\\
44 & 0.2 & 2 & 1.690 & 1.693 & 1.690 & 1.675 (.009) & 1.690 (.007) & 1.690 (.007)\\
44 & 0.4 & 1 & 3.948 & 3.953 & 3.948 & 3.957 (.017) & 3.950 (.011) & 3.931 (.011)\\
44 & 0.4 & 2 & 5.641 & 5.647 & 5.647 & 5.622 (.021) & 5.634 (.012) & 5.644 (.012)\\
    \bottomrule
  \end{tabular}
  \caption{American puts with $K=40$, $r=6\%$ and $50$ exercise dates per year. Bermudan:
    grid of \cref{rem:am-references}; American: smoothed binomial tree with Richardson
    extrapolation. Longstaff--Schwartz: their ``finite difference American'' and LSM columns,
    as printed. Ours: $10^5$ antithetic paths and the four Laguerre basis functions for both
    training and evaluation, standard errors from pair averages; our implementation (NumPy
    \texttt{default\_rng}, seeds $[2001,i,0]$ for training and $[2001,i,1]$ for evaluation in
    row $i=0,\dots,19$).}
  \label{tab:am-ls}
\end{table}

\Cref{fig:am-boundary} looks inside the algorithm for the first row. At $t=0.5$, panel (a)
shows the realised cash flows that step~(2) regresses (minus the exercise value $K-s$, so
that the decision is the sign), the fitted continuation value and the exact one from the
grid. The fit is within a few cents of the truth over the range where paths are, and it
crosses the exercise value close to where the truth does. Panel (b) compares the exercise
boundaries $b(t_k)$, the largest price at which each rule exercises. The LSM boundaries
are noisy and miss the steep rise near maturity, where few paths are close to the
boundary. The value lost is small. Heuristically, payoff and continuation value are equal at
the exact boundary, so exercising on the wrong side of it by a distance $\delta$ loses an
amount of order $\delta^2$ per date. The LSM policy is worth $4.4744$ against the price $4.4778$
(\cref{tab:am-bounds}).

\begin{figure}[ht]
  \centering
  \includegraphics[width=\textwidth]{ch14-boundary}
  \caption{Longstaff--Schwartz put with $S_0=36$, $\sigma=20\%$, $T=1$, $50$ dates; LSM on
    $10^5$ antithetic training paths (seed $[14,0]$). (a) At $t_k=0.5$: $1500$ randomly
    chosen in-the-money realised cash flows (in time-$t_k$ units) minus $K-s$, the fitted
    continuation value with the four Laguerre functions and the exact one, both minus $K-s$;
    exercise where the curve is below zero. (b) Exercise boundary: the critical price from
    the grid, and the largest price at which the LSM rule exercises, with the Laguerre basis
    and with $1,x,x^2$, $x=s/K$.}
  \label{fig:am-boundary}
\end{figure}

\begin{example}[Control variates, and the test case]\label{exa:am-test-case}
The out-of-sample estimator averages independent copies of $Z_{\hat\tau}$, so every method
of \cref{ch:08-antithetic-control} applies to it. The discounted European payoff
$X=e^{-rT}g(S_T)$ has a known mean, the Black--Scholes price, and is computed along the
same paths; we estimate its coefficient on the training paths, a pilot sample independent of
the evaluation paths, so that the controlled estimator stays unbiased for $\hat L$ and
\cref{prop:am-lsm}(i) is not disturbed (\cref{prop:cv-estimated-beta}). On the first row of
\cref{tab:am-ls} with plain paths (seed $3$) the correlation is $0.56$ and the variance
ratio $0.68$: deep in the money, many paths are exercised early and their cash flow is
weakly related to $S_T$. On the test-case put (seed $1$, $10^5$ paths, the Laguerre basis
in $x=S/100$) the correlation is $0.93$ and the variance ratio $0.13$: the control
estimate is $8.3024$ with standard error $0.0110$, against $8.2766$ and $0.0308$ without it.
Both are well below the price $8.34941$, and this is not noise: the fitted policy is worth
$8.3079$ (computed on the grid) and exercises $35\%$ of the paths before maturity, where
the optimal rule exercises essentially never (\cref{rem:am-degenerate}(b)). The early
exercises happen in the last quarter of the year at prices around $80$ (between $76$ and
$88$ for $90\%$ of them), where the time value of the put, measured by the European put, is
between $2$ and $19$ cents: comparable with the error of the fit, which then exercises paths
whose continuation value it underestimates. The policy loses $0.04$, some $10^4$ times the
early-exercise premium, and the in-sample estimate, $8.2729$, does not reveal it. For the test-case call the same setup gives $7.4949$
in sample and $7.4237$ ($0.0330$) out of sample, around the policy value $7.4462$ and below
the Bermudan price $7.4567$.
\end{example}

\demo{14-american-mc.ipynb}{the regression of step~(2) at any date, the in-sample and
out-of-sample estimates for any number of paths and basis, and the exercise boundaries.}

\section{Duality: upper bounds}\label{sec:am-duality}

A lower bound alone does not tell how far below the price it is. Upper bounds come from a
dual problem. Its simplest case is perfect foresight: $\max_kZ_k\ge Z_\tau$ for every
$\tau$, so $\E[\max_kZ_k]\ge U_0$. This bound is poor because the holder who sees the future
exercises each path at its best date. Subtracting a martingale with mean zero before taking
the maximum removes part of that foresight without changing the inequality.

\begin{theorem}[Duality]\label{thm:am-duality}
Let $Z$ be adapted with $\E[\max_k|Z_k|]<\infty$, $U$ its Snell envelope, and $M^*$ the
martingale of \eqref{eq:am-doob}.
\begin{enumerate}[label=(\roman*)]
  \item (\emph{Upper bounds}.) For every martingale $M$ with $M_0=0$, and more generally every
    supermartingale with $M_0=0$,
    \[
      U_0\le\E\Big[\max_{0\le k\le m}(Z_k-M_k)\Big].
    \]
  \item (\emph{The bound is attained, without variance}.)
    $\max_k(Z_k-M^*_k)=U_0$ almost surely. Hence
    $U_0=\min_M\E[\max_k(Z_k-M_k)]$ over martingales with $M_0=0$, attained at $M^*$.
  \item (\emph{The gap}.) If $\E[Z_k^2]<\infty$ for every $k$ and $M$ is a square-integrable
    martingale with $M_0=0$, then
    \[
      0\le\E\Big[\max_k(Z_k-M_k)\Big]-U_0\le\E\Big[\max_k|M_k-M^*_k|\Big]
      \le2\,\big\|M_m-M^*_m\big\|_{L^2}.
    \]
\end{enumerate}
\end{theorem}

\begin{proof}
(i) Every $M_k$ is integrable, so $|\max_k(Z_k-M_k)|\le\sum_k(|Z_k|+|M_k|)$ is integrable. For
$\tau\in\mathcal T_0$, \cref{lem:am-optional} gives $\E[M_\tau]\le M_0=0$, so
$\E[Z_\tau]\le\E[Z_\tau-M_\tau]\le\E[\max_k(Z_k-M_k)]$. Take the supremum over $\tau$.

(ii) By \cref{thm:am-dp}(v), $Z_k-M^*_k\le U_k-M^*_k=U_0-A^*_k\le U_0$ for every $k$, so
$\max_k(Z_k-M^*_k)\le U_0$ surely; by (i) its mean is at least $U_0$. A random variable
bounded above by $U_0$ with mean at least $U_0$ equals $U_0$ almost surely.

(iii) $Z_k-M_k=(Z_k-M^*_k)+(M^*_k-M_k)\le U_0+\max_j|M^*_j-M_j|$ by (ii). Under the
hypothesis, $U$ and hence $M^*$ are square-integrable (by the proof of \cref{thm:am-dp}(i),
$|U_k|\le\E[\max_j|Z_j|\mid\F_k]$), so $N=M^*-M$ is a square-integrable martingale, and by
the Cauchy--Schwarz inequality and Doob's $L^2$ maximal inequality
\citep[Section~4.4]{durrett}, $\E[\max_k|N_k|]\le\E[\max_kN_k^2]^{1/2}\le2\,\E[N_m^2]^{1/2}$.
\end{proof}

The theorem is due to \textcite{rogers-duality}, in continuous time, and
\textcite{haugh-kogan}, in discrete time. Rogers states it over martingales $M$ with $M_0=0$ and
$\E[\sup_t|M_t|]<\infty$, with $\sup_{t\le T}(Z_t-M_t)$ in place of the maximum, and Haugh and
Kogan over supermartingales, with the bound written $\pi_0+\E[\max_k(Z_k-\pi_k)]$. Part (ii)
is the striking part: the optimal martingale $M^*$ is the martingale part of the discounted
option price, the gains of the seller's hedging portfolio, and with it every single path
returns the price. This suggests that a martingale close to $M^*$ gives an upper bound with
small variance, which makes the dual estimator cheap where it works. Part (iii) says the gap is controlled by how well $M$ approximates $M^*$ in
mean square, which is how martingales are built in practice: from an approximation of the
value process \citep{haugh-kogan}, or from an exercise policy \citep{andersen-broadie}.

\paragraph{Martingales in closed form.} $M=0$ gives perfect foresight. Another martingale
with $M_0=0$ that needs no simulation is the discounted European price minus its initial
value, $M_k=e^{-rt_k}P(t_k,S_{t_k})-P(0,S_0)$ with $P(t,s)$ the Black--Scholes price of the
European claim with payoff $g$ at $T$; for a put it removes the part of the foresight that a
European put already prices. On the first row of \cref{tab:am-ls},
\cref{tab:am-bounds} shows that perfect foresight gives $7.78$ for a price of $4.48$, and the
European martingale $4.732$: better, but $5.7\%$ high. \exref{ex:am-rogers} improves it.

\begin{figure}[t]
  \centering
  \includegraphics[width=\textwidth]{ch14-duality}
  \caption{Empirical distribution functions of the dual samples $\max_k(Z_k-M_k)$ for the
    setting of \cref{tab:am-bounds}: the European put martingale ($2000$ outer paths), the
    Andersen--Broadie martingale of the LSM policy estimated with $100$ inner paths ($400$
    outer paths), the same martingale with exact continuation values, and the Doob
    martingale of the grid value function ($2000$ outer paths). Dotted: the price $U_0$.
    Dashed in (b): the value $L_0=4.4744$ of the policy. The nested sizes are kept small to
    keep the figure cheap; \cref{tab:am-bounds} reports a run with $500$ inner paths.}
  \label{fig:am-duality}
\end{figure}

\begin{proposition}[Andersen--Broadie]\label{prop:am-andersen-broadie}
Let $Z$ be adapted with $\E[\max_k|Z_k|]<\infty$. Let a policy be given by events
$A_k\in\F_k$ (exercise at $t_k$ on $A_k$), $k<m$, with $A_m=\Omega$, and let
$\tau_k=\min\{j\ge k:\omega\in A_j\}$. Let
\[
  L_k=\E[Z_{\tau_k}\mid\F_k],\qquad\bar L_k=\E[Z_{\tau_{k+1}}\mid\F_k]\ (k<m),\qquad
  M_0=0,\quad M_{k+1}=M_k+L_{k+1}-\bar L_k .
\]
\begin{enumerate}[label=(\roman*)]
  \item $L_k=Z_k$ on $A_k$ and $L_k=\bar L_k$ off $A_k$; $M$ is a martingale; and with
    $H=\E[\max_k(Z_k-M_k)]$,
    \[
      L_0\le U_0\le H .
    \]
  \item For every $k$, $Z_k-M_k=L_0+(Z_k-L_k)+\sum_{j<k}\1_{A_j}(\bar L_j-Z_j)$. Hence
    $\max_k(Z_k-M_k)\ge L_0$ surely, and
    \[
      0\le H-L_0\le\E\Big[\sum_{k<m}\Big(\1_{A_k^c}(Z_k-\bar L_k)^++\1_{A_k}(\bar L_k-Z_k)^+\Big)\Big].
    \]
    In particular, if $A_k=\{Z_k\ge\bar L_k\}$ for every $k<m$, then $L_0=U_0=H$.
  \item (\emph{Nested simulation is biased high}.) Let $\hat L_k$ ($k<m$) be random variables
    with $\E[|\hat L_k|]<\infty$ and $\E[\hat L_k\mid\F_T]=\bar L_k$, and let $\hat M$ be
    defined like $M$ with $\hat L_k$ in place of $\bar L_k$, both in the increments and in
    $L_k$ off $A_k$. Then
    \[
      \E\Big[\max_k(Z_k-\hat M_k)\Big]\ \ge\ H,\qquad
      \E\Big[\max_k(Z_k-\hat M_k)\Big]-H\le\sum_{k<m}\E\big[|\hat L_k-\bar L_k|\big].
    \]
  \item (\emph{A confidence interval}.) Let $\hat V^{\rm low}$ be the average of $n_1$
    independent copies of $Z_{\tau_0}$ and $\hat V^{\rm up}$ the average of $n_2$ independent
    copies of $\max_k(Z_k-\hat M_k)$ (outer paths with their own inner simulations), with
    standard errors $\hat\sigma_1/\sqrt{n_1}$ and $\hat\sigma_2/\sqrt{n_2}$, and both variances
    finite and positive. Then, as $n_1,n_2\to\infty$, the interval
    $\big[\hat V^{\rm low}-z_\delta\hat\sigma_1/\sqrt{n_1},\ \hat V^{\rm up}+z_\delta\hat\sigma_2/\sqrt{n_2}\big]$
    contains $U_0$ with probability at least $1-\delta+o(1)$.
\end{enumerate}
\end{proposition}

\begin{proof}
(i) On $A_k$, $\tau_k=k$; off $A_k$, $\tau_k=\tau_{k+1}$; both are $\F_k$-measurable events,
which gives the formula for $L_k$. By the tower property
$\E[L_{k+1}\mid\F_k]=\E[Z_{\tau_{k+1}}\mid\F_k]=\bar L_k$, so $M$ has increments with
conditional mean zero, and it is integrable because
$|L_k|,|\bar L_k|\le\E[\max_j|Z_j|\mid\F_k]$. Then $L_0=\E[Z_{\tau_0}]\le U_0\le H$ by
\cref{thm:am-dp}(ii) and \cref{thm:am-duality}(i).

(ii) By (i), $\bar L_j=L_j+\1_{A_j}(\bar L_j-Z_j)$, so
$M_k=\sum_{j<k}(L_{j+1}-L_j)-\sum_{j<k}\1_{A_j}(\bar L_j-Z_j)=L_k-L_0-\sum_{j<k}\1_{A_j}(\bar L_j-Z_j)$,
which is the identity. At $k=\tau_0$ the last two terms vanish ($Z=L$ there, and no
$A_j$ occurs before), so the maximum is at least $L_0$. For the bound, $Z_k-L_k$ is $0$ on
$A_k$ and $Z_k-\bar L_k$ off it, and for real numbers
$\max_k\big(a_k+\sum_{j<k}b_j\big)\le\sum_ka_k^++\sum_jb_j^+$. If $A_k=\{Z_k\ge\bar L_k\}$
every term of the bound vanishes, so $H\le L_0$, and (i) gives equality throughout.

(iii) Off $A_{j+1}$ the increment of $\hat M$ uses $\hat L_{j+1}$ for $L_{j+1}$. With
$\delta_j=\hat L_j-\bar L_j$ for $j<m$ and $\delta_m=0$, the difference
$\varepsilon_k=\hat M_k-M_k$ is
\[
  \varepsilon_k=\sum_{j<k}\big(\1_{A^c_{j+1}}\delta_{j+1}-\delta_j\big)
  =\1_{A^c_k}\delta_k-\delta_0-\sum_{0<j<k}\1_{A_j}\delta_j\qquad(k\ge1),
\]
and $\E[\varepsilon_k\mid\F_T]=0$ because the events $A_j$ are in $\F_T$. The function
$x\mapsto\max_k(Z_k-M_k-x_k)$ is convex, so by the conditional Jensen inequality
$\E[\max_k(Z_k-\hat M_k)\mid\F_T]\ge\max_k(Z_k-M_k)$; take expectations. For the second
claim, $|\max_k(Z_k-\hat M_k)-\max_k(Z_k-M_k)|\le\max_k|\varepsilon_k|\le\sum_{j<m}|\delta_j|$,
since each $\delta_j$ appears at most once in $\varepsilon_k$.

(iv) The two estimators are independent, $\E[\hat V^{\rm low}]=L_0\le U_0$ and
$\E[\hat V^{\rm up}]\ge H\ge U_0$ by (i) and (iii). If $U_0$ is below the interval, then
$\hat V^{\rm low}-z_\delta\hat\sigma_1/\sqrt{n_1}>L_0$, an event whose probability tends to
$\delta/2$ by \cref{thm:mc-clt}(ii); likewise above. The union of the two events has
probability at most $\delta+o(1)$.
\end{proof}

The policy of interest is the LSM rule, with $A_k$ the exercise region of \eqref{eq:am-rule}
and the training paths held fixed (conditionally on $\mathcal G$, as in
\cref{prop:am-lsm}). Then $\bar L_k$ is the continuation value of the LSM policy at the
simulated state, an expectation that the regression did not compute. The algorithm of
\textcite{andersen-broadie} estimates it, along each \emph{outer} path and at each date,
by the mean of $n_{\rm in}$ \emph{inner} paths started from the current state and stopped
by the policy; these estimates are conditionally unbiased, so (iii) applies. If $\hat L_k$ is the mean of
$n_{\rm in}$ inner samples with conditional variance $v_k$, then
$\E|\hat L_k-\bar L_k|\le(\E[v_k]/n_{\rm in})^{1/2}$, and the bias is at most
$\sum_k(\E[v_k])^{1/2}/\sqrt{n_{\rm in}}$, of order $m/\sqrt{n_{\rm in}}$. The cost is of order $n_{\rm out}\,n_{\rm
in}\,m^2/2$ simulated steps. Part~(ii) explains the quality of the bound: the gap is made of
the policy's own mistakes, continuing where its continuation value is below the payoff or
exercising where it is above, measured against its own continuation values. A
self-consistent policy has no gap, and a policy that errs rarely and by little has a small one.

\Cref{tab:am-bounds,fig:am-duality} put the pieces together for the put with $S_0=36$. The
LSM policy of \cref{fig:am-boundary} is worth $4.4744$, $0.0034$ below the price. Its
Andersen--Broadie martingale, with continuation values computed exactly on the grid, gives
the upper bound $4.4786$ with a standard error of $0.0002$ from $2000$ outer paths: the
duality gap of the policy is $0.0042$, and $75\%$ of the outer paths return exactly
$L_0$, which part~(ii) predicts for every path on which the policy makes no mistake. Nested
simulation is another matter: with $100$ inner paths per date the estimator is $4.705$, and
with $500$ it is $4.537$. The inner noise enters through the maximum over $50$ dates and
dominates the gap; reducing it takes many inner paths or variance reduction inside the inner
simulation, at a cost that grows like $m^2$. We keep the nested runs small, with $400$ and
$200$ outer paths, so that tests and figures stay cheap; the standard errors are those of
these runs. With the out-of-sample LSM estimate and the nested estimator with $500$ inner paths,
part~(iv) gives the $95\%$ interval $[4.461,4.558]$, which contains the price; with exact
continuation values it would be $[4.461,4.479]$. The Doob martingale of the grid value
function, finally, illustrates \cref{thm:am-duality}(ii): its $2000$ paths return
$4.47781$ with a standard deviation of $1.6\times10^{-5}$, the error of the grid and of the
linear interpolation.

\begin{table}[ht]
  \centering
  \small
  \begin{tabular}{@{}llccc@{}}
    \toprule
    & & estimate & s.e. & paths\\
    \midrule
    price & grid (Bermudan, $50$ dates) & $4.4778$ & & \\
    \midrule
    lower & LSM policy, exact value (grid) & $4.4744$ & & \\
     & LSM policy, out of sample & $4.4793$ & $0.0094$ & $10^5$\\
    \midrule
    upper & $M=0$ (perfect foresight) & $7.779$ & $0.010$ & $10^5$\\
     & European put martingale & $4.7322$ & $0.0014$ & $10^5$\\
     & Andersen--Broadie, $100$ inner paths & $4.705$ & $0.017$ & $400$\\
     & Andersen--Broadie, $500$ inner paths & $4.537$ & $0.011$ & $200$\\
     & Andersen--Broadie, exact continuation values & $4.4786$ & $0.0002$ & $2000$\\
     & Doob martingale of the grid value function & $4.47781$ & $<10^{-6}$ & $2000$\\
    \bottomrule
  \end{tabular}
  \caption{Lower and upper bounds for the Longstaff--Schwartz put with $S_0=36$,
    $\sigma=20\%$, $T=1$, $50$ dates. The LSM policy is the one of \cref{fig:am-boundary}.
    The two closed-form bounds use $10^5$ paths of their own; the other upper bounds share
    $2000$ outer paths (the nested estimators use the first $400$ or $200$). Our implementation (NumPy \texttt{default\_rng}, seeds $[14,j]$; see
    \texttt{duality\_study}).}
  \label{tab:am-bounds}
\end{table}


\demo{14-american-mc.ipynb}{upper bounds from perfect foresight, the European martingale,
the Andersen--Broadie martingale with nested simulation and with exact continuation values,
and the Doob martingale, for any strike and volatility.}

\furtherreading
\textcite{glasserman-mc}, Chapter~8, is the reference for this chapter: the problem in
Section~8.1, regression in Section~8.6, duality in Section~8.7. \textcite{longstaff-schwartz}
introduced the algorithm used here, after \textcite{carriere} and alongside
\textcite{tsitsiklis-vanroy}, who regress values instead of cash flows.
\textcite{clement-lamberton-protter} prove its convergence and \textcite{glasserman-yu}
study the balance of paths and basis functions. Duality is due to \textcite{rogers-duality}
and \textcite{haugh-kogan}, and the primal--dual algorithm to \textcite{andersen-broadie}.
\textcite{peskir-shiryaev}, Chapter~I, and \textcite{follmer-schied}, Chapter~6, treat
optimal stopping and American claims in general.
